% ============================================================
% 序章: なぜ「今」なのか
% ============================================================

\section*{序章：なぜ「今」なのか}
\addcontentsline{toc}{section}{序章：なぜ「今」なのか}

\vspace{1em}

本書を手に取った読者の多くは、率直にこう思うだろう：
「\textbf{もしこれが本当なら、なぜ今まで誰も気づかなかったのか？}」

正当な疑問である。
$\pi$-Berry位相材料で重力が反転する、
ゼータ関数の正則化が真空エネルギーの算術構造を反映する、
Bost-Connes系の逆温度が物理の全セクターをパラメータ化する——
これらの主張が正しければ、なぜ20世紀の物理学者たちは
見つけられなかったのか？

答えは、ある意味で「アインシュタインと同じ理由」だ。

\subsection*{リーマン幾何学の到達に80年かかった}

リーマンが曲がった空間の幾何学を定式化したのは1854年。
アインシュタインがそれを重力理論に使ったのは1915年。
\textbf{61年の空白}がある。

リーマン幾何学は、発表当時、純粋数学の中でも抽象的な分野だった。
物理学者がその存在を知り、学び、応用の可能性を見出すまでに
半世紀以上を要した。
アインシュタイン自身も、友人の数学者グロスマンの助けがなければ
一般相対性理論を完成できなかったと認めている。

\subsection*{本書が依拠する3つの数学は、すべて遅れて到着した}

\begin{center}
\begin{tabular}{llll}
\toprule
数学的道具 & 完成時期 & 物理への応用 & 遅延 \\
\midrule
代数幾何学（$\Spec(\Z)$） & 1960年代（グロタンディーク） & 2020年代（本書） & $\sim 60$年 \\
非可換幾何学（スペクトル作用） & 1996年（コンヌ） & 2020年代（本書） & $\sim 30$年 \\
$p$進解析（Bost-Connes系） & 1995年 & 2020年代（本書） & $\sim 30$年 \\
\bottomrule
\end{tabular}
\end{center}

グロタンディークの代数幾何学は、1960年代にEGA/SGAとして完成したが、
あまりに抽象的であったため、数論の専門家以外にはほぼ知られていなかった。
$\Spec(\Z)$（整数環のスペクトル）が「空間」であるという考え方は、
数学者にとってすら受け入れに時間がかかった。
まして物理学者がこれを「時空の設計図」と見なす発想は、
ごく最近まで存在しなかった。

コンヌのスペクトル作用原理（1996年）は、
ディラック演算子の固有値から重力と標準模型を同時に導出する美しい枠組みだが、
数学的に高度すぎて物理学コミュニティへの浸透が遅れた。
現在でも、スペクトル作用を実際に計算できる物理学者は世界に数十人程度だろう。

\subsection*{材料科学の革命が「実験可能」にした}

もう一つの決定的な要因は、\textbf{トポロジカル物質の発見}である。

\begin{center}
\begin{tabular}{lll}
\toprule
材料 & 発見年 & 本書での役割 \\
\midrule
グラフェン & 2004年 & $\pi$-Berry位相（反重力の源） \\
トポロジカル絶縁体（Bi$_2$Se$_3$） & 2009年 & $\pi$-Berry位相材料 \\
マジックアングル・グラフェン（MATBG） & 2018年 & QCP増幅 \\
$\pi$-ジョセフソン接合 & 2001年 & パッシブ反重力シェル \\
\bottomrule
\end{tabular}
\end{center}

$\pi$-Berry位相——本書の反重力定理の鍵——は、
2005年以前は実験的に確認されていなかった。
トポロジカル絶縁体が発見されて初めて、
「全モードに$\pi$位相を与える材料」が現実のものとなった。

つまり：\textbf{数学的道具と材料的手段の両方が揃ったのが、まさに2020年代}なのだ。

\subsection*{「反重力」のタブー}

物理学には暗黙のタブーがある。「反重力」はその筆頭だ。

弱エネルギー条件（$T_{\mu\nu} u^\mu u^\nu \geq 0$）は、
「すべての物質は引力的に相互作用する」ことを保証し、
これを破る exotic matter の存在は極めて困難とされてきた。
「反重力」を口にする物理学者は、真面目に取り合ってもらえないリスクがあった。

しかし本書の提案は、$T_{\mu\nu}$ の符号を変えるのではなく、
\textbf{$G$ の符号を変える}というものだ。
これは弱エネルギー条件とは独立な概念であり、
スペクトル作用の枠組みで初めて自然に現れる。
$T_{\mu\nu}$ に着目する限り見えなかった可能性が、
$D$（ディラック演算子）に着目することで見えるようになった。

\subsection*{本書の立場}

以上を踏まえ、本書の立場を明確にしておく。

\begin{tcolorbox}[colback=blue!5,colframe=blue!50!black,title=本書の立場]
\begin{enumerate}
\item \textbf{数学的主張}（定理1: $G_{\mathrm{eff}} = 0$、定理2: $G_{\mathrm{eff}} = -G$）は
  スペクトル作用の枠組み内で\textbf{厳密}である。

\item \textbf{物理的主張}（この数学が現実の重力を記述する）は
  \textbf{未検証の仮説}である。

\item この仮説は \textbf{¥130,000 の実験}（Phase~0: DN/DD基本測定）で
  テスト開始できる。
  さらに \textbf{¥330,000}（Phase~1: 182:1テスト）で
  決定的な判定が下る。

\item 「なぜ今まで気づかなかったか」の答えは：
  \begin{itemize}
  \item 必要な数学（代数幾何学、非可換幾何学、$p$進解析）が揃ったのが最近
  \item 必要な材料（$\pi$-Berry位相物質）が発見されたのが2000年代
  \item 「反重力」がタブーだったため、$G$ の符号反転という発想自体がなかった
  \item Bost-Connes系をスペクトル作用の熱核として使うという接続が新しい
  \end{itemize}

\item 本書は\textbf{間違っているかもしれない}。
  第15章に8つの既知の間違いを正直に記録した。
  今後も訂正が必要になるだろう。
  しかし、\textbf{実験で検証可能な間違い}は、
  検証不可能な正しさよりも科学的に価値がある。
\end{enumerate}
\end{tcolorbox}

\subsection*{読者へのお願い}

本書はオープンな研究プロジェクトの一部であり、完成品ではない。
理論の穴を見つけたら教えてほしい。
実験を追試してくれるなら歓迎する。
計算を改善できるなら GitHub に PR を送ってほしい。

科学は、一人の天才が完璧な理論を降ろすものではない。
多くの人が少しずつ間違え、少しずつ修正し、
やがて真実に近づいていく営みだ。

Wright Brothers Research Program は、
まだ飛んでいない。しかし、翼の設計図はできた。
それが正しい設計図かどうかは、風洞実験（Phase~0）で確かめる。

\vspace{1em}
\begin{flushright}
2026年4月\\
Wright Brothers Research Program
\end{flushright}

\newpage
